\section{本章小结}
\begin{conclusion}
我们从抖动的读数出发，先区分过程变化、观测误差与系统性偏置，再看平滑带来的延迟。接下来把位置和速度组成状态，用 F/Q 推进估计，用 H/R 与新观测比较，均值和协方差一起递推。KF、EKF、UKF 的区别主要在怎样传播统计量，而不是谁的名字更长谁就更好。

你应能核对状态单位、矩阵尺寸、创新和协方差，跑通手算与同一序列的比较。遇到非线性，先写清 h 和雅可比求值位置，再检查 sigma 点是否重建原来的统计量。一般推导可以继续查文献，但曲线是否跟得上、偏置是否残留、缺测时有没有重复使用旧信息，这些不能留到“以后再说”。
\end{conclusion}

\begin{problemset}
\item\label{prob:v2-filter-noise} 未建模加速度、传感器随机误差和固定零偏分别怎样处理？解释为什么三者不能一股脑用“大一点的R”解决。
\item\label{prob:v2-filter-response} 一阶低通初值0、输入恒为2、系数0.25，前两步输出是多少？若采样间隔加倍而系数不变，以秒计的相同响应会怎样变化？高通互补输出各是多少？
\item\label{prob:v2-filter-gain} 先验方差2、观测方差3、H=1，先验值4、观测值9，计算增益、更新状态与后验方差。若H变为1行2列，为什么不能用普通逆将位置观测直接变成完整状态？
\item\label{prob:v2-filter-related} 两个独立无偏估计的方差都为4，平均后的方差是多少？若是同一个估计复制两次，方差是多少？说明交叉协方差怎样进入计算。
\item\label{prob:v2-filter-units} 状态位置用m、速度用m/s，观测直接填mm，会造成什么？写出P的三种不同单位，并说明R=0.04与观测标准差的关系。
\item\label{prob:v2-filter-q} 对本章每段随机加速度模型，采样间隔减半时，Q的位置项、交叉项、速度项分别缩小到原来的多少？说明为何这不能直接用于另一种连续白噪声模型。
\item\label{prob:v2-filter-nonlinear} EKF 与 UKF 分别怎样处理非线性？对斜距观测说明EKF创新为什么用 h 而不是 Hx；再解释UT匹配前两阶统计量为何不保证任意分布的非线性方差准确。
\item\label{prob:v2-filter-missing} 第五帧没有新观测时，估计器应怎样递推？为什么不能拿第四帧重复更新？同时说明协方差变小并不能单独证明跟踪误差变小。
\item\label{prob:v2-filter-scalar} 运行一步更新，把R从0.9改为9，提交两次K、x、P和手算核对；随后尝试R=0，记录程序怎样拒绝无效输入。
\item\label{prob:v2-filter-experiment} 运行linear默认、Q=0、固定0.5 m偏置与range变体，分别保存实际数组和指标。指出加速段与偏置段的误差，解释哪一组更适合当前任务，不能只凭曲线更平下结论。
\end{problemset}

\section{本章验收}
请检查者看你实际保存的输入和输出，并现场改一个参数重跑；图上的文字要对应这次运行，不能拿旧图代替新结果。
\begin{enumerate}
\item 在本章项目目录运行平滑实验，展示原始、低通、尾部均值及误差曲线，改alpha后核对无噪声阶跃90\%响应次数，并演示高通与低通相加还原输入。
\item 解释位置/速度状态、F/H和Q/R的尺寸、单位与生成条件，现场核对scalar的一步预测/更新，再运行linear序列；提交真实RMSE、创新/NIS、协方差最小特征值及Q=0或缺测对照。
\item 用非对角P运行sigma重建与仿射传播，核对解析均值和协方差；展示均匀输入非线性例的近似偏差，指出权重、Cholesky因子取列和正定条件。
\item 在同一range序列上比较EKF/UKF，说明雅可比的求值位置、h的创新和UKF加Q后重新取点的原因。展示一次alpha变体与固定偏置对照；按模型、误差和残差解释结果，不要求某个算法必然获胜。
\end{enumerate}
